\documentclass[a4paper]{article}

\usepackage{graphicx}
\usepackage{amsmath,amssymb}
\usepackage{minted}
\usepackage{multirow}
\usepackage{float}
\usepackage{algorithm}
\usepackage{algpseudocode}
\usepackage{todonotes}
\usepackage{forest}
\usepackage{xstring}
\usepackage{xcolor}
\usepackage[most]{tcolorbox}
\usepackage{xparse}
\usepackage{natbib}
\usepackage{adjustbox}

\usetikzlibrary{graphs,graphdrawing}
\usegdlibrary{layered}

% for external graphviz
\input{/home/donkarlo/Dropbox/repo/nd_language_ai_project/src/nd_language_ai/formal/computer/programming/tex/latex/code_clips/external_graphviz.tex}

% for tags
\input{/home/donkarlo/Dropbox/repo/nd_language_ai_project/src/nd_language_ai/formal/computer/programming/tex/latex/part/control_sequence/kind/semtag/semtag.tex}

% for links
\input{/home/donkarlo/Dropbox/repo/nd_language_ai_project/src/nd_language_ai/formal/computer/programming/tex/latex/code_clips/seehere.tex}

% for nested lists and sections
\input{/home/donkarlo/Dropbox/repo/nd_language_ai_project/src/nd_language_ai/formal/computer/programming/tex/latex/code_clips/deep_structure.tex}

\title{Goal}
\date{}
\author{Mohammad Rahmani}

\begin{document}
   \maketitle

   \section{Definition}
      A goal is a representation of a desired future state or outcome. In planning, the goal constrains which future states count as successful and therefore guides the selection of actions or plans that can reach the desired state \cite{mcdermott-1999-using-regression-match-graphs-to-control-search-in-planning}.

   \section{Functional separation}
      In this project, a goal is treated as a cognitive representation, while the origin and execution of goals are kept in their corresponding biological and control/planning layers.

      \paragraph{Biological and homeostatic origin of goals}
      Physiological needs, survival drives, homeostatic deviations, and allostatic regulation can generate or modulate goals. These mechanisms belong to the life layer rather than to the abstract cognitive representation of a goal.

      \seeherefooter{/home/donkarlo/Dropbox/repo/nd_robotic_ai_learning_project/src/robot/life/life_goals/life_goal.tex}

      \paragraph{Motor goal}
      A motor goal is a desired motor outcome used to organize motor planning and control.

      \seeherefooter{/nd_robotic_ai_learning_project/src/robot/part/mind/cognition/part/object_level/process/part/thinking/planning/motion/motor_goal/motor_goal.tex}

      \paragraph{Goal seeking and navigation}
      Methods that transform a spatial goal and environmental constraints into motion commands belong to motion/navigation planning. Artificial potential fields are therefore documented under path planning rather than in this goal-representation file.

      \seeherefooter{/nd_robotic_ai_learning_project/src/robot/part/mind/cognition/part/object_level/process/part/thinking/planning/motion/navigation/path_planning/kind/artificial_potential_field/artificial_potential_field.tex}

   \bibliographystyle{plainnat}
   \bibliography{/home/donkarlo/Dropbox/repo/nd_shared_research/bibliography/bibliography.bib}
\end{document}
